
\chapter{paths.py}
This file contains no algorithm. It records where the data sits on disk and who the thirteen writers are, and every other module asks it rather than working it out independently.


\paragraph{Why it is a separate file}
Earlier in the project each module computed its own paths by walking upwards from its own location in the directory tree. That works until a file is moved to a different folder, at which point the chain of parent directories is wrong and the module can no longer find the data. Collecting the knowledge in one place means there is exactly one line to change if the data moves, and a module that moves folder does not break.

It is also deliberately not merged into one of the functional modules. Every module needs to know where the weights and profiles live, and if that knowledge lived inside, say, the recogniser, then the gantry driver would have to import the entire recogniser in order to find a directory name.


\section{Data}
\begin{lstlisting}[style=sig]
Data = Path(__file__).resolve().parent.parent / 'CNN'
\end{lstlisting}

\noindent\textbf{Output.} The folder holding the datasets, the trained weights and the style profiles.

Derived from the location of this file rather than written as an absolute path, so the project can be copied to another machine or another drive without editing anything.


\begin{quote}\small\itshape This path is one directory level up from where it was before the modules were flattened into a single folder, and the derivation was adjusted accordingly. A path derived by counting parent directories is correct only for the depth it was written for, which is a known fragility of this approach.\end{quote}

\section{Weights}
\begin{lstlisting}[style=sig]
Weights = Data / 'weights'
\end{lstlisting}

\noindent\textbf{Output.} The sub-folder holding trained checkpoints as \texttt{.npz} files.

A checkpoint is a saved copy of every weight in a trained network. The \texttt{.npz} format is numpy's own archive of named arrays, which is used here because the networks are built from numpy arrays directly and no framework's format is involved.


\section{Profiles}
\begin{lstlisting}[style=sig]
Profiles = Data / 'NOGIT' / 'StyleProfiles10'
\end{lstlisting}

\noindent\textbf{Output.} The folder holding one JSON style profile per writer.

A style profile is the complete description of one person's handwriting produced by style.py: their stored letterforms plus the measurements describing how they lay writing out. JSON is used because the contents are nested lists of numbers and text, and because a profile can then be inspected in a text editor when something looks wrong.


\section{PersonalPages}
\begin{lstlisting}[style=sig]
PersonalPages = Data / 'PersonalDataset'
\end{lstlisting}

\noindent\textbf{Output.} The folder holding the photographed pages and their transcripts.

The pages written specifically for this project, as opposed to the public database. Each writer has a sub-folder containing the photographs and a matching transcript for each one.


\section{Cache}
\begin{lstlisting}[style=sig]
Cache = Data / 'NOGIT'
\end{lstlisting}

\noindent\textbf{Output.} The folder holding caches and other large generated files.

Extracting line crops from pages, and cutting glyphs out of those crops, are the slowest operations in the pipeline. Their results are cached here so that a change to a later stage can be tested without repeating them. Keeping the caches out of the code repository also keeps the repository small.


\section{DatabaseWriters}
\begin{lstlisting}[style=sig]
DatabaseWriters = ['150', '151', '152', '153', '384', '551', '552', '588']
\end{lstlisting}

\noindent\textbf{Output.} The identifiers of the eight writers taken from the public handwriting database.

These are writers from a published corpus of scanned handwriting \cite{marti2002}, identified by number. They provide variety that the project could not collect by hand: different ages, different nationalities, and in particular several strongly joined-up hands, which turn out to be the hardest case for every stage of the pipeline.


\section{PersonalWriters}
\begin{lstlisting}[style=sig]
PersonalWriters = ['abhinav', 'dylan', 'owen', 'thiya', 'yeukita']
\end{lstlisting}

\noindent\textbf{Output.} The names of the five writers photographed for this project.

Volunteers who each filled several pages, photographed with a phone camera. These are the writers the project can test end-to-end, because their pages exist as photographs rather than as clean scans and therefore exercise the whole pipeline from uneven lighting onwards.


\section{AllWriters}
\begin{lstlisting}[style=sig]
AllWriters = DatabaseWriters + PersonalWriters
\end{lstlisting}

\noindent\textbf{Output.} All thirteen writers, database then photographed.

Thirteen in total, which is three more than requirement R6 asks for. The order is fixed so that every table in every results document lists writers in the same sequence.


\section{ValFraction}
\begin{lstlisting}[style=sig]
ValFraction = 0.15
\end{lstlisting}

\noindent\textbf{Output.} The share of data held back from training, 0.15.


\paragraph{What a holdout is, and why the fraction lives here}
Measuring a network on the same data it was trained on says almost nothing, because a network with enough capacity can memorise its training examples. Some data is therefore withheld from training entirely and used only for measurement; that is the holdout, or validation set.

The fraction is defined here, once, rather than in each script that needs a split. If two scripts chose their own fractions, or their own random seeds, the same line could be held out by one and trained on by the other, and every accuracy figure in the project would be quietly wrong.


\section{SplitSeed}
\begin{lstlisting}[style=sig]
SplitSeed = 0
\end{lstlisting}

\noindent\textbf{Output.} The fixed random seed used when splitting data, 0.


\paragraph{Why a random choice is made repeatable}
Deciding which lines are held out is a random choice, but it must be the \emph{same} random choice every time the program runs. A random number generator started from a fixed seed produces the same sequence every time, so the split is both arbitrary, which is what makes it fair, and reproducible, which is what makes it trustworthy.

This also means a holdout set can be recovered later even if it was not saved: replaying the split with the same seed reproduces it exactly.


\section{PersonalWriterPages}
\begin{lstlisting}[style=sig]
def PersonalWriterPages(writer)
\end{lstlisting}

\noindent\textbf{Inputs}
\begin{description}[leftmargin=2.2em,style=nextline,font=\normalfont\ttfamily]\setlength\itemsep{1pt}
  \item[{writer}] \normalfont the lower case name of one photographed writer, for example "owen"
\end{description}

\noindent\textbf{Output.} A generator yielding \texttt{(imagePath, labelPath)} for each of that writer's pages, in sorted filename order.


\paragraph{What it is for}
Finds one writer's pages and pairs each with its transcript. Every stage that needs real handwriting with known text starts here: training, profile building, and the measurement of segmentation accuracy.


\paragraph{Why the order matters}
Several measurement scripts pair the n-th detected line crop on a page with the n-th line of that page's transcript. If the pages were returned in whatever order the file system happened to supply, that pairing would still appear to work but would sometimes be comparing a line against the wrong label, and the resulting accuracy would be wrong in a way that no error message would reveal. Sorting makes the order a property of the filenames rather than of the file system.


\paragraph{Why a generator}
A generator produces its values one at a time as they are asked for, rather than building the whole list first. Here that means a caller can begin processing the first page immediately, and can stop after one page without the cost of listing them all.

